% CLAUDE REPORT::ONETAPE-0001::2026-10-28T22:21Z
% !TEX program = xelatex
% Compiles with XeLaTeX or pdfLaTeX (ASCII source).
\documentclass[11pt]{article}
\usepackage[margin=1in]{geometry}
\usepackage{amsmath,amssymb,amsthm}
\usepackage{xurl}
\newtheorem{theorem}{Theorem}
\newtheorem{proposition}[theorem]{Proposition}
\newtheorem{corollary}[theorem]{Corollary}
\theoremstyle{remark}
\newtheorem*{remark}{Remark}
\newcommand{\SAT}{\mathrm{SAT}}
\newcommand{\DTISP}{\mathrm{DTISP}}
\newcommand{\status}[1]{\textbf{[#1]}}

\title{CLAUDE ONETAPE-0001\\[2pt]
\large Review and extension of the one-tape crossing-sequence notes:\\
nondeterminism, portfolios in sequential access, and the quadratic wall}
\author{CLAUDE (independent verification), at the project owner's request}
\date{2026-10-28 15:21 America/Los\_Angeles}

\begin{document}
\maketitle

\noindent\textbf{Status.} P versus NP remains OPEN. Additive report on two notes submitted by the project
owner (preserved in \url{research/agents/claude/onetape/}). No schedules, allocations, ownership or root
\texttt{main.tex} changed. Not an LSPV run.

\section{Verdict on the notes}

\begin{itemize}
\item Note 1 (one-tape bound $\Omega(N^2/\log N)$ for 3-SAT on the equality family $F_{a,b}$):
\status{PROVED/EXACT}. The splicing argument, the injectivity of crossing sequences, and the union bound
over $L=\Theta(k\log k)$ middle boundaries are correct, and the ``same input'' issue is handled correctly.
\status{COMPUTATIONAL} \texttt{verify\_family.py}: in the exact-width-3 encoding, all $4^k$ inputs have one
length and identical block boundaries, and $F_{a,b}\in\SAT\iff a=b$ for all 1,364 pairs with $k\le5$;
$N_k/(k\log_2k)$ stays bounded (65.0, 44.7, 37.9, 36.6, 34.1 for $k=10,\dots,10^5$); the union bound holds
from $k=28,41,62$ for $q=2^4,2^{10},2^{20}$ states.
\item Note 2 (matching $O(kN)$ upper bound on the family; padding and repetition do not amplify; a single
cut carries at most $O(N)$ bits): \status{PROVED/EXACT}, all statements correct.
\item Literature. The result is a known type. Crossing sequences and the $\Omega(n^2)$ one-tape bound for
palindromes are Hennie's; Cobham proved $T\cdot S=\Omega(n^2)$ for palindromes on off-line machines; and
Santhanam (IPL 2001) reduced palindromes to SAT to get $T\cdot S=\Omega(n^2/\mathrm{polylog}\,n)$ for SAT on
multitape Turing machines (as cited in Williams, CCC 2005). Note 1 says it claims no novelty; that is right.
\end{itemize}

\section{Extension 1: the argument is blind to nondeterminism}

\begin{theorem}\label{thm:nd} \status{PROVED/EXACT}
Every \emph{nondeterministic} one-tape machine deciding 3-SAT in this encoding runs for
$\Omega(N^2/\log N)$ steps on some satisfiable input of length $N_k$, for every large $k$.
\end{theorem}
\begin{proof}
Note 1's proof, with one accepting computation fixed for each $F_{a,a}$. If two fixed computations, for
$a\ne b$, have equal crossing sequences at a middle boundary, splicing them gives an accepting computation
on $F_{a,b}$, which is unsatisfiable. Injectivity, counting and the union bound are unchanged.
\end{proof}

\begin{proposition}\label{prop:ntime} \status{PROVED/EXACT (standard simulations)}
A nondeterministic one-tape machine decides 3-SAT in time $O(N^2\log^2N)$.
\end{proposition}
\begin{proof}[Sketch]
A nondeterministic multitape machine guesses an assignment, sorts the literal occurrences by variable id
(merge sort, $O(N\log N)$), annotates each with its value in one sweep, sorts back by position, and checks
every clause: time $O(N\log N)$. The Hartmanis--Stearns simulation of a $t$-step multitape machine by one
tape costs $O(t^2)$ and applies verbatim to nondeterministic machines.
\end{proof}

\noindent\textbf{Consequence.} For one-tape nondeterministic machines, the complexity of SAT is
$N^{2\pm o(1)}$, and crossing-sequence (fooling-set) arguments are tight there. Any argument of this kind
proves the same bound for nondeterministic machines, so it cannot separate deterministic from
nondeterministic one-tape time for SAT by more than polylogarithmic factors. It measures the cost of moving
information along the tape, not the cost of search.

\section{Extension 2: portfolios in sequential access}

\noindent\textbf{Model.} An off-line machine: a two-way read-only input tape whose head moves one cell per
step, a constant number of work tapes of total space $S\ge\log N$, and optionally a two-way read-only
\emph{guess} tape holding $g$ bits ($\exists^{g}$, the portfolio form used in TPORT-0001). Deterministic, or
nondeterministic with choices made on the fly.

\begin{theorem}\label{thm:seq} \status{PROVED/EXACT}
In this model, any machine deciding 3-SAT in this encoding with time $T$ and space $S$ satisfies
\[
T\cdot S=\Omega\!\left(\frac{N^2}{\log N}\right)
\]
on the inputs $F_{a,a}$, provided the guess has $g\le k/2=\Omega(N/\log N)$ bits (any $g$ for on-the-fly
nondeterminism).
\end{theorem}
\begin{proof}
Define the crossing sequence at an input-tape boundary $j$ as the sequence of full configurations other
than the input-head position (state, work-tape contents and heads, guess-tape head) at the moments the
input head crosses $j$. There are at most $C=2^{O(S+\log g)}=2^{O(S)}$ such configurations. Splicing two
computations with equal crossing sequences at $j$ and the same guess string gives a valid computation on
the spliced input, exactly as on one tape.
For each $a$ fix an accepting computation on $F_{a,a}$ and its guess. Some guess string $y$ is used by a set
$\mathcal A$ of at least $2^{k-g}\ge2^{k/2}$ values of $a$. Within $\mathcal A$, crossing sequences at every
middle boundary are injective in $a$. With $h=\lfloor k/(4\log_2C)\rfloor$, at most $C^{h+1}\le C\,2^{k/4}$
members of $\mathcal A$ have a sequence of length $\le h$ at a given boundary, so by the union bound over the
$L$ middle boundaries some $a\in\mathcal A$ has all of them longer than $h$ once $L\,C\,2^{k/4}<2^{k/2}$,
which holds for large $k$ when $S=o(k)$ (otherwise the claimed bound is trivial). Then
$T\ge L\,h=\Omega(k\log k\cdot k/S)$, i.e.\ $T\cdot S=\Omega(k^2\log k)=\Omega(N^2/\log N)$.
For on-the-fly nondeterminism there is no guess tape and the grouping step is not needed.
\end{proof}

\begin{corollary}[T-PORT holds in sequential access] \status{PROVED/EXACT}
In the sequential-access model, every SAT portfolio of at most $2^{k/2}=2^{\Omega(N/\log N)}$ programs with
space $N^{o(1)}$, in particular the UWI portfolio of about $N$ programs, needs time $N^{2-o(1)}$; with space
$\mathrm{polylog}\,N$, time $N^2/\mathrm{polylog}\,N$.
\end{corollary}

\noindent\textbf{Contrast with TPORT-0001/0002.} With random access to the input, alternation trading with
the portfolio rules provably cannot give the UWI portfolio any exponent above $1$ (ledger L34). With
sequential access the same statement holds up to exponent $2$, by an argument that ignores the portfolio
altogether. The difference is entirely the input-access model: crossing sequences have no counterpart when
any input bit can be read in one step. This is also why the result does not transfer to the ARC/FLMV
allocation, whose model has a random-access (indexed) input.

\section{Extension 3: the quadratic wall, three ways}

\begin{proposition}\label{prop:wall} \status{PROVED/EXACT}
\begin{enumerate}
\item \emph{Counting cap.} For one-tape machines over alphabet $\Sigma$ with $q$ states, the bound obtained by
summing $\log_q(\text{number of inputs a cut must distinguish})$ over cuts is at most
$\sum_j\min(j,N-j)\log_q|\Sigma|=O(N^2)$.
\item \emph{Nondeterminism cap.} By Theorem~\ref{thm:nd} and Proposition~\ref{prop:ntime}, any argument that
also applies to nondeterministic one-tape machines is capped at $O(N^2\log^2N)$ for SAT.
\item \emph{Translation.} If SAT is decidable by a deterministic multitape machine in time $t(N)$, it is
decidable by a one-tape machine in time $O(t(N)^2)$ (Hartmanis--Stearns). Hence a one-tape lower bound of
$N^{2+\delta}$ for SAT would give a multitape lower bound of $N^{1+\delta/2-o(1)}$.
\end{enumerate}
\end{proposition}

\noindent Superlinear deterministic time lower bounds for SAT on multitape machines are, as far as I know,
open: the only general superlinear separation, $\mathrm{DTIME}(n)\ne\mathrm{NTIME}(n)$
(Paul--Pippenger--Szemer\'edi--Trotter), gains less than the $\log N$ factor that the reduction of
nondeterministic linear time to SAT loses. So exceeding quadratic one-tape time for SAT is a known hard problem
in disguise, not a matter of a better family.

\section{What this means for Note 2's research target}

Note 2 asks for a family where ``correctness forces expensive computation after information has been
gathered, surviving arbitrary summaries, rearrangement, and reuse.'' By Proposition~\ref{prop:wall}(3), on one
tape that target is equivalent, up to squaring and polylogarithmic factors, to a superlinear lower bound for
SAT on multitape machines. Two pointers:
\begin{itemize}
\item The one place expansion has given superlinear separations is in the \emph{computation} graph, not the
instance: Hopcroft--Paul--Valiant (time versus space) and Paul--Pippenger--Szemer\'edi--Trotter show that
deterministic computations have small separators and segregators, which nondeterminism can exploit; see also
Santhanam, ``On separators, segregators and time versus space'' (CCC 2001). An expander \emph{formula} only
forces information across cuts, which Proposition~\ref{prop:wall}(1) caps at $O(N^2)$.
\item In the random-access model, where the repository's ARC/FLMV and UWI work lives, the available tools are
alternation trading (TPORT-0001) and its barriers.
\end{itemize}

\section{Ledger rows}

\begin{itemize}
\item L38: Theorem~\ref{thm:nd} (deterministic and nondeterministic one-tape, $\Omega(N^2/\log N)$) --- Theorem.
\item L39: Theorem~\ref{thm:seq} and its corollary (sequential-access portfolios) --- Theorem.
\item L40: ``SAT needs one-tape deterministic time $N^{2+\delta}$ for some $\delta>0$'' --- Open. It follows
from P $\ne$ NP, and it would imply a superlinear multitape lower bound
(Proposition~\ref{prop:wall}(3)); ledger category half-win (true if P $\ne$ NP).
\end{itemize}

\section{Literature}

Read in this session: Williams, ``Better time-space lower bounds for SAT and related problems'', CCC 2005,
\url{https://people.csail.mit.edu/rrw/ccc05.pdf} (footnote describing Santhanam's reduction from
PALINDROMES to SAT and the resulting $\Omega(n^2/\mathrm{poly}\log n)$ time--space product on multitape
machines); Maass, ``Quadratic lower bounds for deterministic and nondeterministic one-tape Turing machines'',
STOC 1984, abstract at \url{https://igi-web.tugraz.at/people/Abstracts/Maass:84a}; R.~Santhanam, ``Lower bounds
on the complexity of recognizing SAT by Turing machines'', IPL 2001 (bibliographic record only,
\url{https://portal.mardi4nfdi.de/wiki/Publication:1603493}).
From memory, verify before load-bearing use: Hennie (1965, crossing sequences); Cobham (1966,
time--space for palindromes); Hartmanis--Stearns (1965, multitape to one-tape in $O(t^2)$);
Hopcroft--Paul--Valiant (1977); Paul--Pippenger--Szemer\'edi--Trotter (FOCS 1983); Santhanam (CCC 2001).

\end{document}
% END_OF_VERIFIED_REPORT::CLAUDE::ONETAPE-0001
